\documentclass[12pt,a4paper]{article}
\usepackage[utf8]{inputenc}
\usepackage[T1]{fontenc}
\usepackage{amsmath,amssymb,amsfonts}
\usepackage{graphicx}
\usepackage{booktabs}
\usepackage{hyperref}
\usepackage[margin=2.5cm]{geometry}
\usepackage[numbers]{natbib}

\newcommand{\Lor}{\mathrm{Lor4D}}

\title{Hierarchical Integration as the Structural Root of Identity Basin Deepening\\
  in Causal Set Dimension Selection}

\author{Gang Zhang}
\date{\today}

\begin{document}
\maketitle

\begin{abstract}
In finite causal-set dimension selection, the identity basin of 4D Lorentzian
causal sets ($\Lor$) deepens systematically with poset size $N$.
We investigate the structural mechanism behind this deepening.
In matched-pair experiments controlling for geometric features (effective
dimension, comparable fraction, antichain width), we find that deeper causal
hierarchy---measured by layer count, mean layer gap, and a composite hierarchy
integration index (HII)---correlates with lower combinatorial freedom at
$r\approx -0.83$ (Pearson, $p<10^{-3}$, stable across three nearwall
calibrations and two independent seed replications).
Controlled regression on 544 samples with family fixed effects confirms that
layer count (partial $r=-0.45$, $p=0.00025$) and mean layer gap (partial
$r=-0.42$, $p=0.00025$) retain significance after controlling for width,
comparable fraction, and effective dimension.
Leave-one-confounder-out analysis identifies antichain width as the primary
confounding channel.
A screening-enhancement scan shows that incorporating hierarchy integration
into the second screening layer reduces the crossing weight from
$\zeta_{\mathrm{cross}}\approx 5$ to $\approx -0.1$, demonstrating that
hierarchical depth is not merely explanatory but has screening potential.
These results suggest that basin deepening is not purely a statistical precision
effect but has a structural root: deeper causal integration compresses
combinatorial degrees of freedom, making the $\Lor$ identity template
increasingly rigid as the causal structure grows.
\end{abstract}

\noindent\textbf{Keywords:}
causal sets, hierarchy, combinatorial entropy, historical sedimentation,
basin deepening, structural screening

% ============================================================
\section{Introduction}
% ============================================================

Causal set theory treats spacetime as a locally finite partially ordered set,
and dimension selection can be phrased as a competition among candidate
families~\cite{Surya2019}.
In our data, the identity basin of 4D Lorentzian causal sets deepens
systematically with poset size $N$: the Mahalanobis gap grows from $+0.31$ at
$N=10$ to $1.93\times 10^8$ at $N=1024$, and the effective potential volume
shrinks as $V_{\mathrm{eff}}\propto N^{-1.66}$.

This deepening may be interpreted as a form of structural sedimentation---the
progressive accumulation of geometric identity through causal structure growth.
But the question of \emph{mechanism} remains open: is the deepening purely a
statistical precision effect (the reference manifold sharpens because more
samples reduce variance), or does it have a structural root in the causal
architecture itself?

The present paper addresses this question.
We show that deeper causal hierarchy---more layers, larger mean layer gap,
stronger cross-layer integration---systematically correlates with lower
combinatorial freedom, even after controlling for geometric confounders.
This provides a candidate structural mechanism for basin deepening: the $\Lor$
template becomes increasingly rigid not just because we measure it more precisely,
but because deeper causal integration compresses the combinatorial degrees of
freedom available to competing structures.

% ============================================================
\section{Methods}
% ============================================================

\subsection{Hierarchy indicators}

For each $N$-element poset $G$, we compute:
\begin{itemize}
\item \textbf{Layer count}: the number of levels in the longest chain
  decomposition.
\item \textbf{Mean layer gap}: the average number of layers spanned by edges in
  the transitive reduction (Hasse diagram).
\item \textbf{Adjacent edge fraction}: fraction of Hasse edges connecting
  adjacent layers.
\item \textbf{Long edge fraction}: fraction of Hasse edges spanning $\geq 2$
  layers.
\item \textbf{Hierarchy Integration Index (HII)}: $z(\text{layer\_count}) +
  z(\text{mean\_layer\_gap}) + z(\text{long\_edge\_fraction}) +
  z(\text{reduction\_edge\_density}) + z(\text{interval\_depth\_mean})$,
  where $z(\cdot)$ denotes within-sample $z$-scoring.
\item \textbf{Locality Dominance Index (LDI)}: $z(\text{adjacent\_edge\_fraction})
  - z(\text{long\_edge\_fraction}) - z(\text{mean\_layer\_gap}) -
  z(\text{interval\_long\_edge\_mean})$.
\end{itemize}

\subsection{Three-layer experimental design}

We test the hierarchy--entropy relationship at three levels of increasing
stringency:

\textbf{Layer~1: Matched-pair test.}
Using Lor2D vs MLR survivor matched pairs (controlling for geometric features),
we compute $\Delta\text{HII}$, $\Delta\text{LDI}$, and $\Delta\log H$ for each
pair and test their correlation.

\textbf{Layer~2: Controlled regression.}
On 544 samples across multiple families and $N$ values, we regress $\log H$ on
hierarchy indicators with controls for width, comparable fraction, effective
dimension, $N$, and family fixed effects.

\textbf{Layer~3: Screening enhancement.}
We test whether incorporating HII or LDI into the second screening layer reduces
the crossing weight $\zeta_{\mathrm{cross}}$.

% ============================================================
\section{Results}
% ============================================================

\subsection{Matched-pair correlations}

Three nearwall calibrations and a rescue protocol yield consistent results
(Table~\ref{tab:c_pairs}):

\begin{table}[h]
\centering
\caption{Matched-pair correlations between hierarchy indicators and
$\Delta\log H$ (MLR minus Lor2D). All permutation $p<0.001$.}
\label{tab:c_pairs}
\begin{tabular}{lccc}
\toprule
Calibration & $n_{\mathrm{pairs}}$ & HII$\to\Delta\log H$ & layer\_count$\to\Delta\log H$ \\
\midrule
Original pilot     & 22 & $-0.757$ & $-0.783$ \\
Nearwall expanded  & 35 & $-0.848$ & $-0.835$ \\
Nearwall rescue    & 50 & $-0.839$ & $-0.827$ \\
\bottomrule
\end{tabular}
\end{table}

Independent seed replication (seeds 20260331 and 20260401) yields identical
Tier-2 correlations: HII $\Delta r = -0.836$, layer count $\Delta r = -0.823$.
Leave-$k$-out removal of the 15 most extreme pairs still yields $r\approx -0.73$
to $-0.75$, confirming the signal is not driven by outliers.

\subsection{Controlled regression}

On 544 samples with the strongest control model
(width + comparable fraction + $d_{\mathrm{eff}}$ + $N$ + family fixed effects):

\begin{table}[h]
\centering
\caption{Partial correlations under full controls.}
\label{tab:c_regression}
\begin{tabular}{lcc}
\toprule
Indicator & Partial $r$ & $p$ \\
\midrule
layer\_count     & $-0.450$ & 0.00025 \\
mean\_layer\_gap & $-0.424$ & 0.00025 \\
HII (composite)  & $+0.086$ & 0.049 \\
\bottomrule
\end{tabular}
\end{table}

The composite HII is weakened under full controls (driven by construction
artefacts in the composite weighting), but its constituent components
(layer\_count, mean\_layer\_gap) retain strong negative partial correlations.

\textbf{LOCO analysis.}
Dropping antichain\_width from controls shifts HII from $+0.086$ to $-0.495$
and layer\_count from $-0.450$ to $-0.766$, identifying antichain width as the
primary confounding channel.

\subsection{Screening enhancement}

Incorporating hierarchy indicators into the second screening layer via
$I_{\mathrm{filter2}} = \text{switch\_zscore} + \eta\cdot z(\text{HII})$
and scanning $\eta$:

\begin{table}[h]
\centering
\caption{Screening enhancement: baseline vs best $\zeta_{\mathrm{cross}}$.}
\label{tab:c_enhance}
\begin{tabular}{lcc}
\toprule
$N$ & Baseline $\zeta_{\mathrm{cross}}$ & Best $\zeta_{\mathrm{cross}}$ \\
\midrule
30 & 2.22 & $+0.03$ \\
40 & 2.88 & $-0.52$ \\
44 & 4.81 & $-0.50$ \\
48 & 5.59 & $-0.69$ \\
\bottomrule
\end{tabular}
\end{table}

The crossing weight drops from $\sim 2$--$6$ to $\sim 0$ or negative,
demonstrating that hierarchy integration has genuine screening potential beyond
post-hoc explanation.

% ============================================================
\section{Discussion}
% ============================================================

The three-layer evidence chain supports a structural mechanism for basin
deepening: deeper causal hierarchy compresses combinatorial degrees of freedom.
This is not merely a restatement of ``more constraints = less entropy''; the
matched-pair design controls for geometric features that could trivially explain
the correlation, and the LOCO analysis identifies the specific confounding
channel (antichain width) through which the effect operates.

The screening enhancement result (Layer~3) is particularly significant: it shows
that hierarchy integration is not just an explanatory variable but has the
potential to enter the screening functional itself.
Whether this potential should be realized---i.e., whether HII should be formally
incorporated into $S_{\mathrm{MD}}$---depends on whether the enhancement
replicates at larger sample sizes, which we leave to future work.

\subsection{Connection to historical sedimentation}

In this data set, basin deepening has two components:
(1)~statistical precision (the reference manifold sharpens as $\Sigma\propto
N^{-1}$), and
(2)~structural rigidity (deeper causal integration makes the $\Lor$ template
harder to imitate).
The present paper provides evidence for component~(2).
Together, they explain why basin deepening is not merely a measurement artefact
but a genuine structural process.

% ============================================================
\section{Conclusion}
% ============================================================

We have shown that deeper causal hierarchy correlates with lower combinatorial
freedom at $r\approx -0.83$ in matched-pair designs, $r\approx -0.45$ in
controlled regressions, and produces measurable screening enhancement.
This provides a candidate structural mechanism for the identity basin deepening
observed in finite causal-set identity selection: the $\Lor$ template becomes
increasingly rigid as causal integration deepens, not just because statistical
precision improves.

\begin{thebibliography}{1}
\bibitem{Surya2019}
S.~Surya,
\emph{The causal set approach to quantum gravity},
\emph{Living Rev.\ Relativ.}\ \textbf{22}, 5 (2019),
arXiv:1903.11544.
\end{thebibliography}

\end{document}
